\documentclass[11pt,a4paper]{article}
\usepackage[utf8]{inputenc}
\usepackage[T1]{fontenc}
\usepackage[catalan]{babel}
\usepackage{lmodern}
\usepackage[sfdefault,lining]{FiraSans}
\usepackage[a4paper,top=2.3cm,bottom=2.2cm,left=1.8cm,right=1.8cm,headheight=1.3cm,headsep=0.4cm,footskip=1cm]{geometry}
\usepackage{amsmath,amssymb}
\usepackage{icomma}
\usepackage{graphicx}
\usepackage[table]{xcolor}
\usepackage{tikz}
\usepackage{fontawesome5}
\usepackage{ccicons}
\usepackage{fancyhdr}
\usepackage{lastpage}
\usepackage{tasks}
\usepackage{needspace}
\usepackage[most]{tcolorbox}
\usepackage{array}
\usepackage[hidelinks,pdftitle={Fitxa d'exercicis de la unitat 6: Introducció a l'àlgebra. Equacions de 1r grau. 1r d'ESO},pdfsubject={Matemàtiques 1r ESO},pdfkeywords={llenguatge algebraic, expressions algebraiques, valor numèric, monomis, polinomis, propietat distributiva, equacions de primer grau, parèntesis, denominadors, problemes}]{hyperref}

\definecolor{navy}{HTML}{1B2748}
\definecolor{accent}{HTML}{E8603A}
\definecolor{teal}{HTML}{1F9E90}
\definecolor{soft}{HTML}{EEF1F8}
\definecolor{muted}{HTML}{5D6685}

\setlength{\parindent}{0pt}
\setlength{\parskip}{3pt}
\settasks{label=\alph*), label-format=\bfseries\color{navy}, label-width=1.3em, item-indent=1.8em, column-sep=0.6em, after-item-skip=5pt, before-skip=3pt, after-skip=4pt}

% ---------- capçalera i peu ----------
\newcommand{\llicencia}{\href{https://creativecommons.org/licenses/by-nc-sa/4.0/deed.ca}{\ccbyncsa}\ \ Llicència Creative Commons Reconeixement-NoComercial-CompartirIgual 4.0 Internacional (CC BY-NC-SA 4.0)}
\pagestyle{fancy}
\fancyhf{}
\fancyhead[L]{\raisebox{-0.35\height}{\includegraphics[height=0.95cm]{logo-gva.png}}}
\fancyhead[C]{\small\color{muted}Fitxa d'exercicis · Unitat 6 · 1r d'ESO}
\fancyhead[R]{\raisebox{-0.35\height}{\includegraphics[height=0.95cm]{logo-vermellar.png}}}
\fancyfoot[L]{\scriptsize\color{muted}\llicencia}
\fancyfoot[R]{\small\color{muted}\thepage}
\renewcommand{\headrulewidth}{0.4pt}
\renewcommand{\headrule}{\color{muted!40}\hrule height \headrulewidth}
\fancypagestyle{primera}{\fancyhf{}\renewcommand{\headrulewidth}{0pt}
  \fancyfoot[L]{\scriptsize\color{muted}\llicencia}
  \fancyfoot[R]{\small\color{muted}\thepage}}

% ---------- nivells i exercicis ----------
\newcommand{\estrelles}[1]{{\color{accent}\ifcase#1\or\faStar\faStar[regular]\faStar[regular]\or\faStar\faStar\faStar[regular]\or\faStar\faStar\faStar\fi}}
\newcounter{exer}
\newcommand{\ex}[1]{\par\needspace{4\baselineskip}\medskip\stepcounter{exer}\noindent{\color{navy}\bfseries\theexer.}\ \ {\small\estrelles{#1}}\ \ }

\newcommand{\seccio}[2]{\par\needspace{10\baselineskip}\bigskip
  \begin{tcolorbox}[colback=navy,colframe=navy,coltext=white,arc=2mm,boxrule=0pt,left=3mm,right=3mm,top=1.2mm,bottom=1.2mm]
  {\large\bfseries #1\quad #2}
  \end{tcolorbox}\vspace{-2pt}}

\newcolumntype{C}[1]{>{\centering\arraybackslash}p{#1}}

% quadre de recordatori (només a la fitxa d'exercicis)
\newcommand{\recorda}[1]{\par\smallskip\begin{tcolorbox}[colback=accent!7,colframe=accent!55,arc=2mm,boxrule=0.5pt,left=2.5mm,right=2.5mm,top=1mm,bottom=1mm]\small{\color{accent}\faLightbulb[regular]}\ \textbf{Recorda:}\ #1\end{tcolorbox}}
\newcommand{\res}[1]{#1}\newcommand{\rest}[1]{#1}

% =====================================================================
% Capçalera FentMates de les fitxes: un color per a cada curs.
% Els colors són els de la pàgina inicial (LEVELS a assets/js/materials.js);
% el logo ve de _eines/logo/pdf/ (python3 _eines/logo/genera.py).
%
% Ús, al final del preàmbul de la unitat:
%   % =====================================================================
% Capçalera FentMates de les fitxes: un color per a cada curs.
% Els colors són els de la pàgina inicial (LEVELS a assets/js/materials.js);
% el logo ve de _eines/logo/pdf/ (python3 _eines/logo/genera.py).
%
% Ús, al final del preàmbul de la unitat:
%   % =====================================================================
% Capçalera FentMates de les fitxes: un color per a cada curs.
% Els colors són els de la pàgina inicial (LEVELS a assets/js/materials.js);
% el logo ve de _eines/logo/pdf/ (python3 _eines/logo/genera.py).
%
% Ús, al final del preàmbul de la unitat:
%   \input{../comu/fentmates.tex}
%   \fmcurs{1eso}
% i, al començament del document:
%   \capcalera{Matemàtiques}{1r ESO}{Fitxa d'exercicis · Unitat 2}{Divisibilitat}
% Opcional abans de \input: \def\fmneed{14} (espai mínim abans de cada secció).
% =====================================================================
\makeatletter
\def\fm@logodir{../../logo/pdf}
\providecommand{\fmneed}{10}
\def\fm@cursnom{}

% clau del curs, color i color del text damunt del color
\newcommand{\fm@curs}[3]{\@namedef{fm@color@#1}{#2}\@namedef{fm@text@#1}{#3}}
\fm@curs{1eso}{E6B84A}{navy}
\fm@curs{2eso}{3B82F6}{white}
\fm@curs{3eso}{FF7A50}{white}
\fm@curs{4eso-a}{A78BFA}{white}
\fm@curs{4eso-b}{6D28D9}{white}
\fm@curs{1batx-mi}{9B1C1C}{white}
\fm@curs{1batx-ccss}{EF6B6B}{white}
\fm@curs{2batx-mii}{1E6B47}{white}
\fm@curs{2batx-ccss}{2BB3A3}{white}

\newcommand{\fmcurs}[1]{%
  \@ifundefined{fm@color@#1}{\PackageError{fentmates}{Curs desconegut: #1}{}}{}%
  \edef\fm@tmp{\noexpand\definecolor{curs}{HTML}{\@nameuse{fm@color@#1}}%
               \noexpand\colorlet{curstext}{\@nameuse{fm@text@#1}}}\fm@tmp
  \colorlet{cursfosc}{curs!72!black}%
  \def\fm@logo{\fm@logodir/vertical-#1.pdf}%
  \def\fm@marca{\fm@logodir/marca-#1.pdf}}

% ---------- capçalera de la primera pàgina ----------
% #1 matèria (petit), #2 curs, #3 tipus de document i unitat, #4 títol
\newcommand{\capcalera}[4]{%
  \gdef\fm@cursnom{#2}%
  \par\noindent
  \begin{tcolorbox}[colback=curs,colframe=curs,boxrule=0pt,arc=2.2mm,width=3.5cm,height=2.15cm,
      valign=center,left=3mm,right=2mm,top=1mm,bottom=1mm,nobeforeafter,box align=center]
    \color{curstext}{\footnotesize #1}\par\vspace{1pt}{\fontsize{19}{21}\selectfont\bfseries #2\par}%
  \end{tcolorbox}\hspace{5mm}%
  \begin{minipage}[c]{\dimexpr\linewidth-3.5cm-5mm-2.9cm\relax}
    {\color{cursfosc}\small\bfseries\MakeUppercase{#3}}\par\vspace{2pt}
    {\color{navy}\fontsize{20}{23}\selectfont\bfseries #4\par}
  \end{minipage}\hfill
  \begin{minipage}[c]{2.6cm}\raggedleft\includegraphics[width=2.6cm]{\fm@logo}\end{minipage}\par
  \vspace{2pt}}

% ---------- barres de secció del color del curs ----------
\renewcommand{\seccio}[2]{\par\needspace{\fmneed\baselineskip}\bigskip
  \begin{tcolorbox}[colback=curs,colframe=curs,coltext=curstext,arc=2mm,boxrule=0pt,left=3mm,right=3mm,top=1.2mm,bottom=1.2mm]
  {\large\bfseries #1\quad #2}
  \end{tcolorbox}\vspace{-2pt}}

% ---------- estrelles de dificultat del color del curs (a Batxillerat no se'n posen) ----------
\@ifundefined{estrelles}{}{%
  \renewcommand{\estrelles}[1]{{\color{curs}\ifcase#1\or\faStar\faStar[regular]\faStar[regular]\or\faStar\faStar\faStar[regular]\or\faStar\faStar\faStar\fi}}}

% ---------- capçalera i peu de les pàgines ----------
% primera pàgina: logos institucionals al peu
\fancypagestyle{primera}{\fancyhf{}\renewcommand{\headrulewidth}{0pt}
  \fancyfoot[L]{\raisebox{-0.32\height}{\includegraphics[height=0.78cm]{logo-gva.png}}\hspace{5mm}%
    \raisebox{-0.32\height}{\includegraphics[height=0.78cm]{logo-vermellar.png}}}
  \fancyfoot[R]{\scriptsize\color{muted}\href{https://creativecommons.org/licenses/by-nc-sa/4.0/deed.ca}{\ccbyncsa}\ \ CC BY-NC-SA 4.0 · fentmates.org\qquad\small\thepage}}
% la resta: el símbol a l'esquerra i el curs a la dreta, amb el títol al centre com abans
\fancyhead[L]{\raisebox{-0.3\height}{\includegraphics[height=0.62cm]{\fm@marca}}}
\fancyhead[R]{\small\bfseries\color{cursfosc}\fm@cursnom}
\renewcommand{\headrule}{\color{curs!45}\hrule height \headrulewidth}
\makeatother

%   \fmcurs{1eso}
% i, al començament del document:
%   \capcalera{Matemàtiques}{1r ESO}{Fitxa d'exercicis · Unitat 2}{Divisibilitat}
% Opcional abans de \input: \def\fmneed{14} (espai mínim abans de cada secció).
% =====================================================================
\makeatletter
\def\fm@logodir{../../logo/pdf}
\providecommand{\fmneed}{10}
\def\fm@cursnom{}

% clau del curs, color i color del text damunt del color
\newcommand{\fm@curs}[3]{\@namedef{fm@color@#1}{#2}\@namedef{fm@text@#1}{#3}}
\fm@curs{1eso}{E6B84A}{navy}
\fm@curs{2eso}{3B82F6}{white}
\fm@curs{3eso}{FF7A50}{white}
\fm@curs{4eso-a}{A78BFA}{white}
\fm@curs{4eso-b}{6D28D9}{white}
\fm@curs{1batx-mi}{9B1C1C}{white}
\fm@curs{1batx-ccss}{EF6B6B}{white}
\fm@curs{2batx-mii}{1E6B47}{white}
\fm@curs{2batx-ccss}{2BB3A3}{white}

\newcommand{\fmcurs}[1]{%
  \@ifundefined{fm@color@#1}{\PackageError{fentmates}{Curs desconegut: #1}{}}{}%
  \edef\fm@tmp{\noexpand\definecolor{curs}{HTML}{\@nameuse{fm@color@#1}}%
               \noexpand\colorlet{curstext}{\@nameuse{fm@text@#1}}}\fm@tmp
  \colorlet{cursfosc}{curs!72!black}%
  \def\fm@logo{\fm@logodir/vertical-#1.pdf}%
  \def\fm@marca{\fm@logodir/marca-#1.pdf}}

% ---------- capçalera de la primera pàgina ----------
% #1 matèria (petit), #2 curs, #3 tipus de document i unitat, #4 títol
\newcommand{\capcalera}[4]{%
  \gdef\fm@cursnom{#2}%
  \par\noindent
  \begin{tcolorbox}[colback=curs,colframe=curs,boxrule=0pt,arc=2.2mm,width=3.5cm,height=2.15cm,
      valign=center,left=3mm,right=2mm,top=1mm,bottom=1mm,nobeforeafter,box align=center]
    \color{curstext}{\footnotesize #1}\par\vspace{1pt}{\fontsize{19}{21}\selectfont\bfseries #2\par}%
  \end{tcolorbox}\hspace{5mm}%
  \begin{minipage}[c]{\dimexpr\linewidth-3.5cm-5mm-2.9cm\relax}
    {\color{cursfosc}\small\bfseries\MakeUppercase{#3}}\par\vspace{2pt}
    {\color{navy}\fontsize{20}{23}\selectfont\bfseries #4\par}
  \end{minipage}\hfill
  \begin{minipage}[c]{2.6cm}\raggedleft\includegraphics[width=2.6cm]{\fm@logo}\end{minipage}\par
  \vspace{2pt}}

% ---------- barres de secció del color del curs ----------
\renewcommand{\seccio}[2]{\par\needspace{\fmneed\baselineskip}\bigskip
  \begin{tcolorbox}[colback=curs,colframe=curs,coltext=curstext,arc=2mm,boxrule=0pt,left=3mm,right=3mm,top=1.2mm,bottom=1.2mm]
  {\large\bfseries #1\quad #2}
  \end{tcolorbox}\vspace{-2pt}}

% ---------- estrelles de dificultat del color del curs (a Batxillerat no se'n posen) ----------
\@ifundefined{estrelles}{}{%
  \renewcommand{\estrelles}[1]{{\color{curs}\ifcase#1\or\faStar\faStar[regular]\faStar[regular]\or\faStar\faStar\faStar[regular]\or\faStar\faStar\faStar\fi}}}

% ---------- capçalera i peu de les pàgines ----------
% primera pàgina: logos institucionals al peu
\fancypagestyle{primera}{\fancyhf{}\renewcommand{\headrulewidth}{0pt}
  \fancyfoot[L]{\raisebox{-0.32\height}{\includegraphics[height=0.78cm]{logo-gva.png}}\hspace{5mm}%
    \raisebox{-0.32\height}{\includegraphics[height=0.78cm]{logo-vermellar.png}}}
  \fancyfoot[R]{\scriptsize\color{muted}\href{https://creativecommons.org/licenses/by-nc-sa/4.0/deed.ca}{\ccbyncsa}\ \ CC BY-NC-SA 4.0 · fentmates.org\qquad\small\thepage}}
% la resta: el símbol a l'esquerra i el curs a la dreta, amb el títol al centre com abans
\fancyhead[L]{\raisebox{-0.3\height}{\includegraphics[height=0.62cm]{\fm@marca}}}
\fancyhead[R]{\small\bfseries\color{cursfosc}\fm@cursnom}
\renewcommand{\headrule}{\color{curs!45}\hrule height \headrulewidth}
\makeatother

%   \fmcurs{1eso}
% i, al començament del document:
%   \capcalera{Matemàtiques}{1r ESO}{Fitxa d'exercicis · Unitat 2}{Divisibilitat}
% Opcional abans de \input: \def\fmneed{14} (espai mínim abans de cada secció).
% =====================================================================
\makeatletter
\def\fm@logodir{../../logo/pdf}
\providecommand{\fmneed}{10}
\def\fm@cursnom{}

% clau del curs, color i color del text damunt del color
\newcommand{\fm@curs}[3]{\@namedef{fm@color@#1}{#2}\@namedef{fm@text@#1}{#3}}
\fm@curs{1eso}{E6B84A}{navy}
\fm@curs{2eso}{3B82F6}{white}
\fm@curs{3eso}{FF7A50}{white}
\fm@curs{4eso-a}{A78BFA}{white}
\fm@curs{4eso-b}{6D28D9}{white}
\fm@curs{1batx-mi}{9B1C1C}{white}
\fm@curs{1batx-ccss}{EF6B6B}{white}
\fm@curs{2batx-mii}{1E6B47}{white}
\fm@curs{2batx-ccss}{2BB3A3}{white}

\newcommand{\fmcurs}[1]{%
  \@ifundefined{fm@color@#1}{\PackageError{fentmates}{Curs desconegut: #1}{}}{}%
  \edef\fm@tmp{\noexpand\definecolor{curs}{HTML}{\@nameuse{fm@color@#1}}%
               \noexpand\colorlet{curstext}{\@nameuse{fm@text@#1}}}\fm@tmp
  \colorlet{cursfosc}{curs!72!black}%
  \def\fm@logo{\fm@logodir/vertical-#1.pdf}%
  \def\fm@marca{\fm@logodir/marca-#1.pdf}}

% ---------- capçalera de la primera pàgina ----------
% #1 matèria (petit), #2 curs, #3 tipus de document i unitat, #4 títol
\newcommand{\capcalera}[4]{%
  \gdef\fm@cursnom{#2}%
  \par\noindent
  \begin{tcolorbox}[colback=curs,colframe=curs,boxrule=0pt,arc=2.2mm,width=3.5cm,height=2.15cm,
      valign=center,left=3mm,right=2mm,top=1mm,bottom=1mm,nobeforeafter,box align=center]
    \color{curstext}{\footnotesize #1}\par\vspace{1pt}{\fontsize{19}{21}\selectfont\bfseries #2\par}%
  \end{tcolorbox}\hspace{5mm}%
  \begin{minipage}[c]{\dimexpr\linewidth-3.5cm-5mm-2.9cm\relax}
    {\color{cursfosc}\small\bfseries\MakeUppercase{#3}}\par\vspace{2pt}
    {\color{navy}\fontsize{20}{23}\selectfont\bfseries #4\par}
  \end{minipage}\hfill
  \begin{minipage}[c]{2.6cm}\raggedleft\includegraphics[width=2.6cm]{\fm@logo}\end{minipage}\par
  \vspace{2pt}}

% ---------- barres de secció del color del curs ----------
\renewcommand{\seccio}[2]{\par\needspace{\fmneed\baselineskip}\bigskip
  \begin{tcolorbox}[colback=curs,colframe=curs,coltext=curstext,arc=2mm,boxrule=0pt,left=3mm,right=3mm,top=1.2mm,bottom=1.2mm]
  {\large\bfseries #1\quad #2}
  \end{tcolorbox}\vspace{-2pt}}

% ---------- estrelles de dificultat del color del curs (a Batxillerat no se'n posen) ----------
\@ifundefined{estrelles}{}{%
  \renewcommand{\estrelles}[1]{{\color{curs}\ifcase#1\or\faStar\faStar[regular]\faStar[regular]\or\faStar\faStar\faStar[regular]\or\faStar\faStar\faStar\fi}}}

% ---------- capçalera i peu de les pàgines ----------
% primera pàgina: logos institucionals al peu
\fancypagestyle{primera}{\fancyhf{}\renewcommand{\headrulewidth}{0pt}
  \fancyfoot[L]{\raisebox{-0.32\height}{\includegraphics[height=0.78cm]{logo-gva.png}}\hspace{5mm}%
    \raisebox{-0.32\height}{\includegraphics[height=0.78cm]{logo-vermellar.png}}}
  \fancyfoot[R]{\scriptsize\color{muted}\href{https://creativecommons.org/licenses/by-nc-sa/4.0/deed.ca}{\ccbyncsa}\ \ CC BY-NC-SA 4.0 · fentmates.org\qquad\small\thepage}}
% la resta: el símbol a l'esquerra i el curs a la dreta, amb el títol al centre com abans
\fancyhead[L]{\raisebox{-0.3\height}{\includegraphics[height=0.62cm]{\fm@marca}}}
\fancyhead[R]{\small\bfseries\color{cursfosc}\fm@cursnom}
\renewcommand{\headrule}{\color{curs!45}\hrule height \headrulewidth}
\makeatother

\fmcurs{1eso}

\begin{document}
\thispagestyle{primera}

\capcalera{Matemàtiques}{1r ESO}{Fitxa d'exercicis · Unitat 6}{Introducció a l'àlgebra.\\ Equacions de 1r grau}
\vspace{10pt}
Nom i cognoms: \rule{8.4cm}{0.4pt}\hfill Grup: \rule{1.6cm}{0.4pt}\hfill Data: \rule{2.2cm}{0.4pt}
\vspace{8pt}

\begin{tcolorbox}[colback=soft,colframe=soft,arc=2mm,left=3mm,right=3mm,top=1.5mm,bottom=1.5mm]
\small Exercicis graduats per nivells de dificultat:\quad \estrelles{1}\ bàsic \qquad \estrelles{2}\ mitjà \qquad \estrelles{3}\ avançat.\\[2pt] Resol-los al quadern mostrant tots els passos i comprova les solucions de les equacions. Fes els càlculs sense calculadora.
\end{tcolorbox}

% =====================================================
\seccio{1}{Llenguatge algebraic}

\ex{1} Expressa en llenguatge algebraic (fes servir la lletra $x$):\begin{tasks}(2)
  \task El doble d'un nombre
  \task Un nombre més 7
  \task El triple d'un nombre menys 2
  \task La meitat d'un nombre
  \task El quadrat d'un nombre
  \task Un nombre menys la seua quarta part
  \task El consecutiu d'un nombre
  \task El doble d'un nombre més el seu triple
\end{tasks}\recorda{Entre un nombre i una lletra no escrivim el punt de multiplicar: $3\cdot x=3x$. «El doble d'un nombre més 3» és $2x+3$, però «el doble de la suma d'un nombre i 3» és $2(x+3)$.}

\ex{1} Escriu amb paraules què expressa cada expressió algebraica:\begin{tasks}(3)
  \task $x+5$
  \task $3x$
  \task $x-1$
  \task $\dfrac{x}{3}$
  \task $2(x+4)$
  \task $x^2-1$
\end{tasks}

\ex{2} L'edat de la Núria és $x$ anys. Expressa en llenguatge algebraic:\begin{tasks}(2)
  \task L'edat de la Núria d'ací a 6 anys.
  \task L'edat de la Núria fa 4 anys.
  \task L'edat de la seua àvia, que té 5 vegades l'edat de la Núria.
  \task L'edat del seu cosí, que té la meitat d'anys que ella.
  \task La suma de les edats de la Núria i la seua àvia.
  \task L'edat que tindrà l'àvia d'ací a 10 anys.
\end{tasks}

\ex{2} Expressa en llenguatge algebraic:\begin{tasks}(2)
  \task Dos nombres consecutius.
  \task Un nombre parell i el parell següent.
  \task Un nombre senar.
  \task La suma de tres nombres consecutius.
  \task El preu de $x$ quaderns si cada un costa $1,50$\,€.
  \task Els cèntims que hi ha en $x$ monedes de 20 cèntims.
\end{tasks}

\ex{2} Escriu una fórmula per a:\begin{tasks}(1)
  \task Els euros que paguem per $n$ entrades de cine a 8\,€ cadascuna.
  \task El nombre de rodes que hi ha en $c$ cotxes i $m$ motos.
  \task Els minuts que hi ha en $h$ hores.
  \task El perímetre i l'àrea d'un quadrat de costat $l$.
\end{tasks}

\ex{3} Observa la sèrie de figures fetes amb punts:\par\smallskip{\centering \begin{tikzpicture}[baseline=0]\fill[accent] (0,0.2) circle (0.09); \fill[accent] (0,0.6) circle (0.09);\fill[navy] (0.40,0.00) circle (0.09);\fill[navy] (0.40,0.40) circle (0.09);\fill[navy] (0.40,0.80) circle (0.09);\end{tikzpicture}\qquad\qquad \begin{tikzpicture}[baseline=0]\fill[accent] (0,0.2) circle (0.09); \fill[accent] (0,0.6) circle (0.09);\fill[navy] (0.40,0.00) circle (0.09);\fill[navy] (0.40,0.40) circle (0.09);\fill[navy] (0.40,0.80) circle (0.09);\fill[navy] (0.80,0.00) circle (0.09);\fill[navy] (0.80,0.40) circle (0.09);\fill[navy] (0.80,0.80) circle (0.09);\end{tikzpicture}\qquad\qquad \begin{tikzpicture}[baseline=0]\fill[accent] (0,0.2) circle (0.09); \fill[accent] (0,0.6) circle (0.09);\fill[navy] (0.40,0.00) circle (0.09);\fill[navy] (0.40,0.40) circle (0.09);\fill[navy] (0.40,0.80) circle (0.09);\fill[navy] (0.80,0.00) circle (0.09);\fill[navy] (0.80,0.40) circle (0.09);\fill[navy] (0.80,0.80) circle (0.09);\fill[navy] (1.20,0.00) circle (0.09);\fill[navy] (1.20,0.40) circle (0.09);\fill[navy] (1.20,0.80) circle (0.09);\end{tikzpicture}\par}\smallskip a) Quants punts té cada figura? Quants en tindrà la figura 4? \ b) Escriu una expressió per als punts de la figura $n$. \ c) Quants punts tindrà la figura 10? I la figura 50?

\ex{3} Tradueix al llenguatge algebraic:\begin{tasks}(1)
  \task La diferència entre el triple d'un nombre i la seua meitat.
  \task El doble de la suma d'un nombre i el seu consecutiu.
  \task La suma dels quadrats de dos nombres $a$ i $b$.
  \task El quadrat de la suma de dos nombres $a$ i $b$.
\end{tasks}


% =====================================================
\seccio{2}{Expressions algebraiques. Valor numèric}

\ex{1} Calcula el valor numèric de $4x-3$ per a: \quad a) $x=2$ \quad b) $x=0$ \quad c) $x=-1$ \quad d) $x=5$\recorda{Per a calcular el valor numèric, substituïm la lletra pel nombre i fem les operacions respectant la jerarquia. Els nombres negatius, sempre entre parèntesis: si $x=-3$, \ $x^2=(-3)^2=9$.}

\ex{1} Calcula el valor numèric de cada expressió per als valors indicats:\begin{tasks}(2)
  \task $x^{2}+1$ \ per a \ $x=3$
  \task $2a+b$ \ per a \ $a=4,\ b=-3$
  \task $10-3x$ \ per a \ $x=-2$
  \task $\dfrac{x}{2}+5$ \ per a \ $x=8$
  \task $3(x-1)$ \ per a \ $x=4$
  \task $a\cdot b-a$ \ per a \ $a=-2,\ b=5$
\end{tasks}

\ex{2} Copia i completa la taula de valors:\par\smallskip{\centering\renewcommand{\arraystretch}{1.35}\begin{tabular}{|c|C{0.9cm}|C{0.9cm}|C{0.9cm}|C{0.9cm}|C{0.9cm}|C{0.9cm}||}\hline $x$ & $-2$ & $-1$ & $0$ & $1$ & $2$ & $3$\\\hline $2x+1$ &  &  &  &  &  & \\\hline $x^2-x$ &  &  &  &  &  & \\\hline\end{tabular}\par}

\ex{2} Fes servir les fórmules per a calcular:\begin{tasks}(1)
  \task L'àrea d'un triangle, $A=\dfrac{b\cdot h}{2}$, si $b=10$\,cm i $h=7$\,cm.
  \task El perímetre d'un rectangle, $P=2b+2h$, si $b=12$\,cm i $h=4,5$\,cm.
  \task El volum d'un cub, $V=a^3$, si $a=4$\,cm.
\end{tasks}

\ex{2} Una companyia de telèfon cobra 5\,€ al mes més $0,10$\,€ per cada minut de trucada. Escriu l'expressió del que paguem en un mes si parlem $m$ minuts. Quant paguem si parlem 120 minuts? I si en parlem 45?

\ex{3} Quina expressió dona el valor de $y$ en cada taula?\par\smallskip{\centering\renewcommand{\arraystretch}{1.3}\begin{tabular}{|c|C{0.8cm}|C{0.8cm}|C{0.8cm}|C{0.8cm}|}\hline $x$ & 1 & 2 & 3 & 4\\\hline $y$ & 2 & 6 & 10 & 14\\\hline\end{tabular}\qquad\begin{tabular}{|c|C{0.8cm}|C{0.8cm}|C{0.8cm}|C{0.8cm}|}\hline $x$ & 0 & 1 & 2 & 3\\\hline $y$ & 1 & 2 & 5 & 10\\\hline\end{tabular}\par}

\ex{3} Busca provant un valor de $x$ que faça que: \ a) $3x+2$ valga 20; \ b) $x^2$ valga 49; \ c) $2x+5$ i $x+9$ tinguen el mateix valor.


% =====================================================
\seccio{3}{Monomis i polinomis}

\ex{1} Copia i completa la taula:\par\smallskip{\centering\renewcommand{\arraystretch}{1.35}\begin{tabular}{|c|C{2.2cm}|C{2.2cm}|C{1.4cm}|}\hline Monomi & Coeficient & Part literal & Grau\\\hline $6x^3$ & & & \\\hline $-a^2b$ & & & \\\hline $xy$ & & & \\\hline $-7$ & & & \\\hline $\tfrac{3}{4}x^2y^2$ & & & \\\hline $-z$ & & & \\\hline\end{tabular}\par}\recorda{En el monomi $-5x^2y$: el coeficient és $-5$, la part literal és $x^2y$ i el grau és la suma dels exponents, $2+1=3$. Si no hi ha nombre, el coeficient és 1 (o $-1$).}

\ex{1} Quines d'aquestes expressions són monomis?\begin{tasks}(4)
  \task $5x^2$
  \task $x+1$
  \task $-3ab$
  \task $\dfrac{4}{x}$
  \task $8$
  \task $x^2y^3$
  \task $2^x$
  \task $\dfrac{x}{5}$
\end{tasks}

\ex{2} Agrupa els monomis semblants: \quad $4x,\ \ -2x^2,\ \ 3xy,\ \ x,\ \ 5x^2,\ \ -xy,\ \ 7,\ \ -3x$

\ex{2} Escriu: \ a) dos monomis semblants a $-2x^3$; \ b) un monomi de grau 5 amb coeficient 4; \ c) un monomi de grau 0; \ d) un monomi de grau 2 amb dues lletres diferents.

\ex{2} Indica el nombre de termes, el grau i el terme independent de cada polinomi:\begin{tasks}(2)
  \task $3x^2-5x+1$
  \task $7-x^3$
  \task $x^4+2x^2-x$
  \task $5x-8$
\end{tasks}

\ex{3} Escriu un monomi semblant a $6x^2y$ que: \ a) sumat amb ell done 0; \ b) sumat amb ell done $10x^2y$; \ c) restat d'ell done $x^2y$.

\ex{3} Calcula el valor numèric: \ a) del monomi $-3x^2y$ per a $x=-2$, $y=5$; \ b) del polinomi $x^3-2x+1$ per a $x=-1$; \ c) del polinomi $2x^2-3x$ per a $x=\frac{1}{2}$.


% =====================================================
\seccio{4}{Operacions amb monomis. Reduir expressions}

\ex{1} Calcula:\begin{tasks}(4)
  \task $5x+3x$
  \task $9a-4a$
  \task $x+x+x$
  \task $2x-7x$
  \task $4x^{2}+x^{2}$
  \task $-3ab+8ab$
  \task $6y-y$
  \task $-x-x$
\end{tasks}\recorda{Només podem sumar o restar monomis semblants: sumem els coeficients i deixem la part literal. $3x+5x=8x$, però $3x+2$ no es pot reduir.}

\ex{1} Redueix:\begin{tasks}(2)
  \task $3x+2+5x-1$
  \task $7-2x+4-x$
  \task $x^{2}+3x-2x^{2}+x$
  \task $5a+2b-3a+b$
\end{tasks}

\ex{1} Multiplica:\begin{tasks}(3)
  \task $3\cdot 5x$
  \task $2x\cdot 4x$
  \task $(-3a)\cdot 2a$
  \task $x\cdot x^{2}$
  \task $5x\cdot (-2y)$
  \task $(-4x^{2})\cdot (-x)$
\end{tasks}

\ex{2} Aplica la propietat distributiva:\begin{tasks}(3)
  \task $4(x+3)$
  \task $2(5x-1)$
  \task $-3(x+2)$
  \task $-(4x-7)$
  \task $x(x+5)$
  \task $2x(3x-1)$
\end{tasks}\recorda{$a(b+c)=ab+ac$: multipliquem per \emph{cada} terme del parèntesi. Un signe menys davant canvia tots els signes: $-(x-5)=-x+5$.}

\ex{2} Trau parèntesis i redueix:\begin{tasks}(2)
  \task $2(x+4)+3x$
  \task $5x-(2x+3)$
  \task $3(x-2)+2(x+5)$
  \task $4(2x+1)-(x-6)$
\end{tasks}

\ex{2} Troba l'error en cada igualtat i corregeix-la:\begin{tasks}(2)
  \task $3x+4=7x$
  \task $2(x+5)=2x+5$
  \task $-(x-3)=-x-3$
  \task $x\cdot x=2x$
\end{tasks}

\ex{3} Redueix:\begin{tasks}(1)
  \task $3(2x-1)-2(x-4)+x$
  \task $5-2(3x+1)+4(x-2)$
  \task $x(x+2)-3(x-1)$
\end{tasks}

\ex{3} Escriu de manera reduïda el perímetre de: \ a) un rectangle de base $2x+1$ i altura $x-1$; \ b) un triangle de costats $x$, $x+3$ i $2x-1$. Calcula'ls per a $x=4$.

\ex{3} Endevinalla: pensa un nombre, suma-li 5, multiplica el resultat per 2, resta-li 10 i divideix entre 2. Què obtens? Explica per què passa sempre fent servir l'àlgebra.


% =====================================================
\seccio{5}{Equacions. Elements i solució}

\ex{1} Indica els membres, els termes, la incògnita i el grau de cada equació: \quad a) $5x-3=2x+9$ \quad b) $4(a+1)=a-2$\recorda{Els \textbf{membres} són les expressions de cada costat del signe $=$; els \textbf{termes}, cada sumand; la \textbf{incògnita}, la lletra; i la \textbf{solució}, el valor que fa certa la igualtat.}

\ex{1} Comprova si el valor indicat és solució de l'equació:\begin{tasks}(2)
  \task $3x-1=8$, \ $x=3$
  \task $2x+5=x+1$, \ $x=-4$
  \task $5-x=2x$, \ $x=2$
  \task $\dfrac{x}{3}+1=x-3$, \ $x=6$
\end{tasks}

\ex{1} Troba mentalment la solució:\begin{tasks}(3)
  \task $x+4=10$
  \task $x-3=8$
  \task $5x=30$
  \task $\dfrac{x}{2}=7$
  \task $12-x=5$
  \task $2x+1=11$
\end{tasks}

\ex{2} Quin dels valors $-2$, $0$, $1$ o $3$ és la solució de cada equació?\begin{tasks}(2)
  \task $4x-3=2x+3$
  \task $2(x+1)=x$
  \task $7x=x$
  \task $3-x=2x$
\end{tasks}

\ex{2} Digues si cada igualtat és una identitat o una equació:\begin{tasks}(2)
  \task $2(x+1)=2x+2$
  \task $3x+1=3x+2$
  \task $x+x=2x$
  \task $4x-x=12$
\end{tasks}

\ex{3} Escriu una equació de primer grau: \ a) que tinga per solució $x=5$; \ b) que tinga per solució $x=-2$ i un parèntesi; \ c) que no tinga solució.


% =====================================================
\seccio{6}{Equacions equivalents. Resolució}

\ex{1} Resol:\begin{tasks}(4)
  \task $x+8=15$
  \task $x-6=-2$
  \task $x+9=4$
  \task $-5+x=-11$
\end{tasks}\recorda{Transposició: el que \emph{suma} passa \emph{restant} i el que \emph{resta} passa \emph{sumant}; el que \emph{multiplica} passa \emph{dividint} (amb el seu signe) i el que \emph{divideix} passa \emph{multiplicant}.}

\ex{1} Resol:\begin{tasks}(3)
  \task $6x=42$
  \task $-4x=20$
  \task $9x=-27$
  \task $\dfrac{x}{5}=3$
  \task $\dfrac{x}{4}=-3$
  \task $-x=8$
\end{tasks}

\ex{1} Resol:\begin{tasks}(4)
  \task $2x+3=11$
  \task $5x-4=21$
  \task $3x+10=1$
  \task $7-2x=13$
\end{tasks}

\ex{2} Resol i comprova la solució:\begin{tasks}(2)
  \task $5x+2=3x+10$
  \task $8x-5=6x+1$
  \task $4x+7=x-8$
  \task $2x-9=5x+3$
\end{tasks}

\ex{2} Resol:\begin{tasks}(2)
  \task $3x+4-x=x+9$
  \task $6-2x+5=3x+1$
  \task $x+2x+3x=4x+10$
  \task $9x-3-4x=2x+12$
\end{tasks}

\ex{2} Resol (la solució pot ser una fracció: dona-la irreductible):\begin{tasks}(4)
  \task $4x=10$
  \task $6x+1=2x+3$
  \task $3x-2=5$
  \task $10x=-4$
\end{tasks}

\ex{3} Resol. Què passa en les dues últimes?\begin{tasks}(2)
  \task $2x+5-7x=3-x-6$
  \task $12-3x+x=4x-6$
  \task $x+1=x+5$
  \task $3x+6-x=2x+6$
\end{tasks}

\ex{3} Escriu l'equació que representa la balança i troba quant pesa cada capsa $x$ (les boles pesen 1\,kg).\par\smallskip{\centering \begin{tikzpicture}[scale=0.85]\draw[fill=soft,draw=muted] (3,2.2)--(2.6,-0.4)--(3.4,-0.4)--cycle; \draw[line width=2pt] (0.6,2.2)--(5.4,2.2); \fill[accent] (3,2.2) circle (0.09);\draw[muted] (1.20,2.2)--(-0.10,0.5) (1.20,2.2)--(2.50,0.5);\draw[thick,fill=soft] (-0.10,0.5) arc (180:360:1.3 and 0.25);\draw[fill=navy!80,draw=navy] (0.05,0.53) rectangle +(0.38,0.38); \node[text=white,font=\small\itshape] at (0.24,0.72) {x};\draw[fill=navy!80,draw=navy] (0.53,0.53) rectangle +(0.38,0.38); \node[text=white,font=\small\itshape] at (0.72,0.72) {x};\draw[fill=navy!80,draw=navy] (1.01,0.53) rectangle +(0.38,0.38); \node[text=white,font=\small\itshape] at (1.20,0.72) {x};\draw[fill=accent!60,draw=accent] (1.68,0.72) circle (0.18); \node[font=\tiny] at (1.68,0.72) {1};\draw[fill=accent!60,draw=accent] (2.16,0.72) circle (0.18); \node[font=\tiny] at (2.16,0.72) {1};\draw[fill=accent!60,draw=accent] (0.24,1.18) circle (0.18); \node[font=\tiny] at (0.24,1.18) {1};\draw[fill=accent!60,draw=accent] (0.72,1.18) circle (0.18); \node[font=\tiny] at (0.72,1.18) {1};\draw[muted] (4.80,2.2)--(3.50,0.5) (4.80,2.2)--(6.10,0.5);\draw[thick,fill=soft] (3.50,0.5) arc (180:360:1.3 and 0.25);\draw[fill=navy!80,draw=navy] (3.65,0.53) rectangle +(0.38,0.38); \node[text=white,font=\small\itshape] at (3.84,0.72) {x};\draw[fill=accent!60,draw=accent] (4.32,0.72) circle (0.18); \node[font=\tiny] at (4.32,0.72) {1};\draw[fill=accent!60,draw=accent] (4.80,0.72) circle (0.18); \node[font=\tiny] at (4.80,0.72) {1};\draw[fill=accent!60,draw=accent] (5.28,0.72) circle (0.18); \node[font=\tiny] at (5.28,0.72) {1};\draw[fill=accent!60,draw=accent] (5.76,0.72) circle (0.18); \node[font=\tiny] at (5.76,0.72) {1};\draw[fill=accent!60,draw=accent] (3.84,1.18) circle (0.18); \node[font=\tiny] at (3.84,1.18) {1};\draw[fill=accent!60,draw=accent] (4.32,1.18) circle (0.18); \node[font=\tiny] at (4.32,1.18) {1};\draw[fill=accent!60,draw=accent] (4.80,1.18) circle (0.18); \node[font=\tiny] at (4.80,1.18) {1};\draw[fill=accent!60,draw=accent] (5.28,1.18) circle (0.18); \node[font=\tiny] at (5.28,1.18) {1};\draw[fill=accent!60,draw=accent] (5.76,1.18) circle (0.18); \node[font=\tiny] at (5.76,1.18) {1};\draw[fill=accent!60,draw=accent] (3.84,1.64) circle (0.18); \node[font=\tiny] at (3.84,1.64) {1};\end{tikzpicture}\par}


% =====================================================
\seccio{7}{Equacions amb parèntesis}

\ex{1} Resol:\begin{tasks}(2)
  \task $3(x+2)=18$
  \task $5(x-1)=20$
  \task $2(x+7)=x+12$
  \task $-2(x-3)=10$
\end{tasks}\recorda{Primer traiem parèntesis (propietat distributiva) i després resolem com sempre. Compte amb el signe menys: $-2(x-3)=-2x+6$.}

\ex{2} Resol:\begin{tasks}(2)
  \task $4(x-1)=2(x+3)$
  \task $3(2x+1)=5x+7$
  \task $6-(x+2)=x$
  \task $2x-(5-x)=10$
\end{tasks}

\ex{2} Resol:\begin{tasks}(2)
  \task $5(x+2)-3x=16$
  \task $7x-2(x+4)=12$
  \task $3(x-4)+2(x+1)=5$
  \task $4(2x-3)=5(x+3)$
\end{tasks}

\ex{3} Resol. Què passa en l'última equació?\begin{tasks}(1)
  \task $2(3x-1)-3(x+2)=x+4$
  \task $5-2(x-3)=3(x-8)+5$
  \task $4(x+1)-3(x-2)=2(x+5)$
  \task $3(x-1)-(2x+5)=x-8$
\end{tasks}

\ex{3} Resol:\begin{tasks}(1)
  \task $x-2(x-5)=3(4-x)+6$
  \task $6(x+1)-4(2x-1)=0$
\end{tasks}


% =====================================================
\seccio{8}{Equacions amb denominadors}

\ex{1} Resol:\begin{tasks}(4)
  \task $\dfrac{x}{3}=5$
  \task $\dfrac{x}{4}+2=5$
  \task $\dfrac{2x}{3}=8$
  \task $\dfrac{x}{2}-1=4$
\end{tasks}\recorda{Multipliquem \emph{tots} els termes dels dos membres pel m.c.m. dels denominadors (també els que no tenen denominador). Si un numerador té més d'un terme, el posem entre parèntesis.}

\ex{2} Resol:\begin{tasks}(4)
  \task $\dfrac{x}{2}+\dfrac{x}{4}=6$
  \task $\dfrac{x}{3}-\dfrac{x}{6}=2$
  \task $\dfrac{x}{2}+\dfrac{x}{3}=10$
  \task $x-\dfrac{x}{5}=8$
\end{tasks}

\ex{2} Resol:\begin{tasks}(4)
  \task $\dfrac{x+3}{2}=5$
  \task $\dfrac{x-1}{3}=x-5$
  \task $\dfrac{2x+1}{5}=3$
  \task $\dfrac{x}{2}=\dfrac{x+6}{5}$
\end{tasks}

\ex{2} Resol:\begin{tasks}(2)
  \task $\dfrac{x}{2}+3=\dfrac{x}{3}+5$
  \task $\dfrac{x+1}{4}=\dfrac{x-1}{2}$
\end{tasks}

\ex{3} Resol:\begin{tasks}(3)
  \task $\dfrac{x+2}{3}-\dfrac{x-1}{4}=2$
  \task $x-\dfrac{x+2}{3}=6$
  \task $\dfrac{3x-1}{2}-\dfrac{x+5}{4}=\dfrac{x}{4}$
\end{tasks}

\ex{3} Resol:\begin{tasks}(2)
  \task $\dfrac{x}{2}+\dfrac{x}{4}+\dfrac{x}{8}=14$
  \task $\dfrac{2(x-1)}{3}=\dfrac{x+4}{2}$
\end{tasks}


% =====================================================
\seccio{9}{Resolució de problemes amb equacions}

\ex{1} Si a un nombre li sumem 12, obtenim 30. Quin és el nombre?\recorda{Passos: 1) llig i entén l'enunciat; 2) tria la incògnita ($x=\ldots$); 3) planteja l'equació; 4) resol-la; 5) comprova la solució i respon la pregunta.}

\ex{1} El triple d'un nombre menys 8 és 25. Quin és el nombre?

\ex{1} La suma de dos nombres consecutius és 73. Quins són?

\ex{2} La suma de tres nombres consecutius és 159. Quins són?

\ex{2} La Laura té 6 anys més que el seu germà. Entre els dos sumen 30 anys. Quants anys té cadascun?

\ex{2} Un bolígraf i una llibreta costen junts $4,50$\,€. La llibreta costa el doble que el bolígraf. Quant costa cada cosa?

\ex{2} El perímetre d'un rectangle és 46\,cm i la base fa 5\,cm més que l'altura. Quant mesuren els costats?

\ex{2} Repartim 90 cromos entre tres amics. El segon en rep el doble que el primer i el tercer, el triple que el primer. Quants en rep cadascun?

\ex{3} Una mare té 42 anys i el seu fill, 10. D'ací a quants anys la mare tindrà el triple d'edat que el fill?

\ex{3} La meitat d'un nombre més la seua quarta part és 21. Quin és el nombre?

\ex{3} En una classe hi ha 28 alumnes. Si hi hagueren 4 xiques més, hi hauria tantes xiques com xics. Quantes xiques i quants xics hi ha?

\ex{3} En un aparcament hi ha cotxes i motos: en total, 30 vehicles i 96 rodes. Quants cotxes i quantes motos hi ha?

\ex{3} Els angles d'un triangle mesuren $x$, $2x$ i $3x$. Quant mesura cada angle? (Recorda que els tres angles d'un triangle sumen $180^\circ$.)

\vfill
\begin{center}\color{muted}\small\faCheckCircle[regular]\ \ Ànim! Si et costa algun exercici, repassa els apunts interactius de la unitat.\end{center}

\end{document}
